\documentclass[12pt]{article}

\usepackage[T1]{fontenc}
\usepackage[utf8]{inputenc}
\usepackage{amsmath,amssymb,amsthm,enumitem}
\usepackage[margin=1in]{geometry}
\usepackage[colorlinks=true,linkcolor=blue]{hyperref}

\setlength{\parindent}{0pt}
\setlength{\parskip}{5pt}
\setlength{\abovedisplayskip}{7pt}
\setlength{\belowdisplayskip}{7pt}
\setlist{leftmargin=1.5em,itemsep=4pt,topsep=4pt}

\newcommand{\op}{\diamond}
\newcommand{\E}[1]{\mathrm{E#1}}
\newcommand{\fin}{\mathrm{fin}}
\newcommand{\im}{\operatorname{im}}
\newcommand{\Fix}{\operatorname{Fix}}

\theoremstyle{definition}
\newtheorem{problem}{Problem}
\newtheorem*{warning}{Warning}
\newtheorem*{goal}{Final goal}

\title{A Finite Counterexample Challenge for Equation 677}
\author{}
\date{May 2026}

\begin{document}
\maketitle

\begin{problem}[Counterexample question]
Assume that the finite implication
\[
\E{677}\models_{\fin}\E{255}
\]
is false.  Find a finite nonempty magma $(M,\op)$ satisfying
\[
\E{677}:\qquad
x=y\op\bigl(x\op((y\op x)\op y)\bigr)
\]
for all $x,y\in M$, together with a point $x\in M$ for which
\[
\E{255}:\qquad
x=((x\op x)\op x)\op x
\]
fails.  Prefer a smallest example, or a branch analysis that narrows what the smallest
example must look like.
\end{problem}

This is the companion problem to the proof challenge.  Here the working attitude is the
opposite: suppose a bad finite point exists, then force its local diagram until either a
complete Cayley table appears or the branch collapses.

\begin{warning}
The identities below are necessary conditions for a counterexample, not a substitute for
checking the full law $\E{677}$.  Do not assume associativity, commutativity, an identity
element, inverses, right cancellation, or quasigroup structure.  Every cancellation below
is left cancellation, justified by finite left division.
\end{warning}

\section*{Hints}

\begin{enumerate}[label=\textbf{\arabic*.}]
\item \textbf{Start with the universal tools.}
In any finite $\E{677}$-magma, every left translation $L_a(t)=a\op t$ is bijective, and
\[
a\backslash b=b\op((a\op b)\op a). \tag{L}
\]
Thus left cancellation is available.  Also keep the two global identities
\[
z=(y\op z)\op\bigl((y\op(y\op z))\op y\bigr), \tag{T}
\]
\[
(y\op(y\op z))\op y
=z\op\bigl(((y\op z)\op z)\op(y\op z)\bigr). \tag{K}
\]
These are the main algebraic engines for filling a partial table.

\item \textbf{Normalize the bad point.}
Fix a point $x$ at which $\E{255}$ fails, and put
\[
a=x\op x,
\qquad
q=a\op x=((x\op x)\op x),
\qquad
p=x\backslash x.
\]
Then the badness condition is
\[
q\op x\ne x. \tag{B0}
\]
Moreover, if $u\op x=x$, then $u=q$.  Hence, under (B0), there is no solution to
$u\op x=x$ at all.  Equivalent bad conditions are
\[
q\backslash x\ne x,
\qquad
((x\op x)\op q)\op(x\op x)\ne x,
\qquad
(q\op x)\op q\ne p. \tag{B1}
\]
These give three quick tests for any proposed table.

\item \textbf{Build the principal left orbit.}
Define
\[
c_0=x,
\qquad
c_{k+1}=x\op c_k,
\qquad
d_k=c_k\op x.
\]
Let $d$ be the period of the $L_x$-orbit of $x$.  The recurrence
\[
c_{k-1}=c_k\op d_{k+1} \tag{O}
\]
holds for all $k$ modulo $d$, and
\[
p=c_{d-1},
\qquad
q=d_1=c_{d-2},
\qquad
p\op c_1=q. \tag{O1}
\]
In a bad point, $d\ge4$.  Also every orbit edge must avoid the two forbidden values
\[
d_k=c_k\op x\notin\{x,c_{k-1}\}. \tag{O2}
\]

For a labelled search, the first symmetry break is usually
\[
x=0,
\qquad
0\op0=1,
\qquad
0\op1=2,
\qquad
0\op2=3,
\]
with the rest of the $L_x$-cycle chosen according to the period $d$.

\item \textbf{Exploit the defective right column.}
Since no element maps to $x$ under $R_x(t)=t\op x$, the value $x$ is missing from
\[
\im R_x.
\]
Finiteness therefore forces a nontrivial right collision.  Choose either any collision
or the first repetition in the right tail
\[
\rho_0=x,
\qquad
\rho_{n+1}=\rho_n\op x.
\]
Thus there are $y\ne z$ and $e$ with
\[
y\op x=z\op x=e. \tag{C0}
\]
Then (T) gives the rectangle
\[
(y\op e)\op y=(z\op e)\op z=e\backslash x, \tag{C1}
\]
and the splitter condition
\[
y\op e\ne z\op e. \tag{C2}
\]
This is not a contradiction.  It is the shape a genuine counterexample must repeatedly
thread.

\item \textbf{Keep the fixed maps fixed-point-free.}
Define
\[
H(t)=(x\op t)\op x,
\qquad
F(t)=x\op(t\op x).
\]
They satisfy
\[
x\backslash t=t\op H(t),
\qquad
F(L_x(t))=L_x(H(t)). \tag{F}
\]
Thus $F$ and $H$ are conjugate finite self-maps.  At a bad point neither has a fixed
point.  Equivalently, a counterexample must avoid
\[
H(t)=t
\qquad\text{and}\qquad
F(t)=t
\]
for every $t\in M$.

\item \textbf{Use the permutation-label viewpoint.}
Because every left translation is bijective, write
\[
y\op t=\sigma(y)(t),
\qquad
\sigma(y)\in\operatorname{Sym}(M).
\]
The label map $y\mapsto\sigma(y)$ is injective.  A right collision (C0) becomes
\[
\sigma(y)(x)=\sigma(z)(x)=e,
\qquad
\sigma(y)(e)\ne\sigma(z)(e),
\]
and the rectangle becomes
\[
\sigma(\sigma(y)(e))(y)=e\backslash x
=\sigma(\sigma(z)(e))(z). \tag{P}
\]
This is often the cleanest way to search: construct an injective family of permutations
whose pointwise multiplication table satisfies $\E{677}$.

\item \textbf{Period four is the first test arena.}
If $d=4$, write
\[
c_0=x,
\qquad
c_1=a,
\qquad
c_2=q,
\qquad
c_3=p,
\qquad
r=q\op x,
\qquad
s=p\op x.
\]
Every bad period-four candidate must satisfy
\[
x\op a=q,
\quad
a\op x=q,
\quad
a\op r=x,
\quad
x\op q=p,
\quad
x\op p=x,
\quad
p\op a=q, \tag{P4}
\]
and
\[
r\notin\{x,a,q\}. \tag{P4a}
\]
Thus either $r=p$ or $r$ leaves the principal orbit $\{x,a,q,p\}$.

In the delayed branch $r=p$, one also has
\[
q\op a=q,
\qquad
s=(q\op q)\op q=a\op(q\op q),
\qquad
s\notin\{x,a,q,p\}. \tag{P4b}
\]
In the external branch $r\notin\{x,a,q,p\}$, set
\[
h=r\op q,
\qquad
t=q\op a,
\qquad
u=x\op r.
\]
Useful forced cells are
\[
q\op s=a,
\qquad
(p\op q)\op p=s,
\qquad
a\op h=r. \tag{P4c}
\]
In particular, a four-element orbit-only table cannot be bad; the right tail must create
external structure.

\item \textbf{Period five gives compact internal collisions.}
If $d=5$, write
\[
c_0=x,
\qquad
c_1=a,
\qquad
c_2=b,
\qquad
c_3=q,
\qquad
c_4=p,
\]
and put
\[
e=b\op x,
\qquad
g=q\op x,
\qquad
s=p\op x.
\]
Then a bad period-five candidate satisfies
\[
a\op x=q,
\quad
x\op q=p,
\quad
x\op p=x,
\quad
p\op a=q, \tag{P5}
\]
\[
g=a\op((b\op a)\op b),
\qquad
q\op((a\op q)\op a)=x,
\qquad
(a\op q)\op a=x\op(g\op q), \tag{P5a}
\]
and the exclusions
\[
e\notin\{x,a\},
\qquad
g\notin\{x,b\},
\qquad
s\notin\{x,q\}. \tag{P5b}
\]
The orbit recurrence also gives
\[
b\op g=a,
\qquad
q\op s=b,
\qquad
p\op a=q. \tag{P5c}
\]

If the right tail stays inside the principal orbit through $g$, then
$g\in\{a,q,p\}$.  The three internal branches begin as follows:
\[
\begin{array}{rcl}
g=q &\Longrightarrow& b\op q=a,
\quad (a\op q)\op a=(q\op q)\op q=q\backslash x,
\quad a\op q\ne q\op q,\vspace{1mm}\\
g=a &\Longrightarrow& b\op a=a,
\quad (q\op a)\op q=b\op x=a\backslash x,
\quad q\op a\ne b,\vspace{1mm}\\
g=p &\Longrightarrow& b\op p=a,
\quad q\op s=b,
\quad s\notin\{x,q\}.
\end{array} \tag{P5d}
\]
These branches are good places to apply the collision rectangle (C1) by hand.

\item \textbf{For period six and beyond, use uniform ground constraints.}
If $d=6$, write
\[
c_0=x,
\quad
c_1=a,
\quad
c_2=b,
\quad
c_3=c,
\quad
c_4=q,
\quad
c_5=p,
\]
and put
\[
h=c\op x,
\qquad
g=q\op x.
\]
Then
\[
a\op x=q,
\quad
x\op q=p,
\quad
x\op p=x,
\quad
p\op a=q, \tag{P6}
\]
\[
h=a\op((b\op a)\op b),
\qquad
g=b\op((c\op b)\op c), \tag{P6a}
\]
\[
q\op((a\op q)\op a)=x,
\qquad
(a\op q)\op a=x\op(g\op q), \tag{P6b}
\]
with
\[
b\op x\notin\{x,a\},
\quad
h\notin\{x,b\},
\quad
g\notin\{x,c\},
\quad
p\op x\notin\{x,q\}. \tag{P6c}
\]
For longer periods, repeatedly use
\[
d_{k+2}=c_k\op\bigl((c_{k+1}\op c_k)\op c_{k+1}\bigr), \tag{U1}
\]
\[
x=d_k\op\bigl((c_k\op d_k)\op c_k\bigr), \tag{U2}
\]
\[
(c_k\op d_k)\op c_k=x\op\bigl((d_k\op x)\op d_k\bigr). \tag{U3}
\]

\item \textbf{Use the known construction atlas as search pressure.}
The public Equation 677 database should be read as empirical calibration, not as a
proof.  Its current finite examples all satisfy $\E{255}$, and they cluster into a few
recognizable positive families:
\[
\begin{array}{ll}
\text{linear field magmas} & x\op y=(1-\alpha)x+\alpha y,\quad \Phi_{10}(\alpha)=0,\\
\text{Steiner line magmas} & S(2,5,n)\text{ blocks, each a copy of the unique size-5 magma},\\
\text{pencil magmas} & \text{one pivot with a small number of size-5 petals},\\
\text{fiber bundles} & \text{product or twisted product structure with visible fibers},\\
\text{piecewise-linear magmas} & x\op y=x+f(y-x)\text{ with }f\text{ split by residues or cosets},\\
\text{near-field style magmas} & \text{large }L_0,R_0\text{ cycles and non-linear field-extension behavior}.
\end{array}
\]
Most of these examples are right-cancellative, and hence cannot contain a bad point of
the kind being sought here.  The non-right-cancellative examples are still useful: they
show how right columns can collapse while every left row remains bijective, but in every
known case the collapses still thread the collision rectangles and preserve $\E{255}$.

\item \textbf{Colored-slope lift branch.}
One concrete counterexample route is now isolated enough to be worth exhausting by SAT.
Let
\[
M=\mathbb F_p\times\mathbb F_q,
\qquad
(y,i)\op(x,j)=\bigl(Ay+Bx,\,O_{[y:x]}(i,j)\bigr)
\]
for $(y,x)\ne(0,0)$, with a separate operation on the zero fiber.  For the primary
target
\[
p=19,
\quad q=7,
\quad A=7,
\quad B=4,
\quad \lambda_{\mathrm{bad}}=5,
\]
the affine base $y*x=7y+4x$ satisfies $\E{677}$ on $\mathbb F_{19}$, and
\[
1*1=11,
\qquad 11*1=5,
\qquad 5*1=1.
\]
Thus a fourth-power failure is forced if the color table at slope $5$ is
\[
O_5=C_3,
\qquad
O_4(u,v)=u+2v,
\qquad
O_{12}(u,v)=6u+2v,
\]
and the remaining slope operations can be completed with row permutations satisfying all
slope identities.  The seed itself passes the whole $\lambda=5$ identity:
\[
O_5\bigl(i,O_{12}(j,O_4(O_5(i,j),i))\bigr)=j.
\]
For $i=0$ this is
\[
O_5(0,j)=3j,
\quad O_4(3j,0)=3j,
\quad O_{12}(j,3j)=5j,
\quad O_5(0,5j)=j,
\]
and for $i\ne0$ it is
\[
O_5(i,j)=j+i,
\quad O_4(j+i,i)=j+3i,
\quad O_{12}(j,j+3i)=j+6i,
\quad O_5(i,j+6i)=j.
\]
Moreover, for $s\ne0$ the final $C_3$ row avoids $s$, so any full completion gives a
bad point.

The obstruction is not the bad slope but the coupled side block.  The slope maps force
the simultaneous block
\[
\{1,2,4,6,8,12,14,16,17\},
\]
rather than independent repairs.  In particular, from the fixed $O_4$ and $O_{12}$ one
gets the derived constraints
\[
O_{11}(O_1(i,j),i)=2(i-j),
\]
so each column of $O_{11}$ is a permutation and $O_1$ is determined by those columns,
and
\[
O_{18}\bigl(j,O_1(6i+2j,i)\bigr)=4(j-i).
\]
The naive local repair
\[
O_6(u,v)=v,
\qquad O_{11}(u,v)=u+v,
\]
with $O_{14}$ forced by $\lambda=6$, and the analogous $\lambda=17$ repair, satisfies the
two neighboring equations but fails at the next compatibility equation $\lambda=2$ once
$O_1$ is forced.  More generally, the affine/projection-side family
\[
O_{11}(u,v)=au+bv+c,
\qquad a,b\ne0,
\]
together with $O_6(u,v)=v$ and the forced $O_{14}$ is UNSAT for all $6\cdot6\cdot7=252$
cases under the primary seed.  Therefore a completion, if it exists, must make at least
one of the $O_{11},O_6,O_{14}$ block genuinely non-affine or non-projection.

The search file is \texttt{explore\_colored\_magma.py}.  It checks the slope equations,
the full magma identity including the zero fiber, and the final $\E{255}$ failure.  Use
the affine sweep to keep the dead repair family closed, then run deep branches over the
remaining SAT instance.  The most useful branch variable is $O_{11}$: the primary seed
has a homogeneous color-scalar symmetry, so the branch $O_{11}(0,0)=0$ plus the normalized
nonzero branch $O_{11}(0,0)=1$ covers the same quotient as all nonzero values.  For
machine exhaustion, split an entire column of $O_{11}$ into permutations and shard it
across workers; this is a partition of the remaining search, not an extra mathematical
assumption.

\item \textbf{Use minimality as pressure.}
A smallest counterexample cannot be commutative, right-cancellative, a quasigroup, or a
loop, because the bad right column $R_x$ misses $x$.  The bad point cannot be idempotent,
and its $L_x$-period is at least $4$.

As computational calibration only: the public database export has $841$ known finite
$\E{677}$ magmas, all marked as satisfying $\E{255}$.  Sizes $2,3,4,6,8,10$ have been
proved empty for $\E{677}$, while sizes $12,14,15$ remain open or only partially ruled
out.  Bad-witness searches have found no finite counterexample through order $9$.  At
order $10$, periods $4,5,6,8,9$ are closed in the existing search logs, leaving the broad
period-$7$ and period-$10$ branches as the interesting order-$10$ frontier.  These facts
should guide effort, but a proposed counterexample still needs an independent table
check.

\item \textbf{What counts as a finished counterexample.}
A completed answer should provide a full Cayley table for $M$ and verify:
\[
\text{$\E{677}$ holds for all pairs }(x,y),
\qquad
\text{$\E{255}$ fails at the chosen point }x.
\]
It should also report the data
\[
a=x\op x,
\quad
q=a\op x,
\quad
p=x\backslash x,
\quad
q\op x,
\quad
\text{the $L_x$-period }d,
\]
and exhibit at least one right collision fiber satisfying the two displayed collision constraints.
Also record the fingerprints: whether any pair-generated submagmas have size $5$; whether
the table is translation-invariant of the form $u\op v=u+f(v-u)$ after relabelling;
whether the right-column collisions respect visible fibers; and the cycle structures of
$L_x$ and $R_x$ at the bad point.  If a branch is claimed impossible, the contradiction
should use only the constraints above and left cancellation.
\end{enumerate}

\begin{goal}
Construct a finite table satisfying all of $\E{677}$ and failing $\E{255}$, or push one
of the named branches to contradiction.  A true counterexample must pass every local
test above; a failed branch should explain exactly which forced equations collide.
\end{goal}

\end{document}
